% !TEX root = ../main.tex
\chapter{Kernel Abstraction Layer}
\label{ch:kernel_abstraction}

\chref{ch:op_registry} described the op-level surface: \(\sim 260\)
operations, each carrying a schema, dispatched through the backend
to either \MLX\ or \Accelerate. The dispatch hides which library
runs an op, but it does not hide the structural similarities
between ops in the same category. Every unary elementwise op
(\code{relu}, \code{exp}, \code{sigmoid}, \ldots) shares the same
control structure: iterate over elements, apply a scalar function,
write to the output. Every reduction op (\code{sum}, \code{mean},
\code{max}, \ldots) shares an axis-iteration pattern. Every binary
op shares a broadcasting machinery.

If each op's C++ file replicated this control structure, the engine
would contain dozens of nearly identical loops. The \emph{kernel
abstraction layer} extracts the structure into a small set of
templated kernels --- one per arity class --- and reduces each
per-op file to a declaration of its scalar function plus any
op-specific knowledge.

This chapter describes the kernel layer. It is the layer below
\code{ops/} and above \code{backend/} in the DAG of
\chref{ch:layered_dag}: ops call into kernel templates; kernel
templates call into backend primitives. \secref{sec:kernel_abstraction:why}
motivates the layer; \secref{sec:kernel_abstraction:ikernel}
describes the abstract interface;
\secref{sec:kernel_abstraction:unary}
through~\ref{sec:kernel_abstraction:variadic} treat each kernel
class; \secref{sec:kernel_abstraction:contig} treats the
contiguity inference that the kernels consult; and
\secref{sec:kernel_abstraction:primitives} treats the
data-dependent primitives that do not fit the standard kernel
templates.

\section{Why a Kernel Layer}
\label{sec:kernel_abstraction:why}

The argument for a kernel abstraction is straightforward
duplication-elimination, but the consequences are large:

\begin{itemize}
  \item \textbf{Code volume.} Without the layer, the \(\sim 60\)
    unary ops would each contain a full elementwise loop with
    stride handling, broadcasting, dtype dispatch, and AMP-aware
    casting. With the layer, each op file declares its scalar
    function (\code{[](float x) \{ return x > 0 ? x : 0; \}} for
    ReLU) and the template handles the loop. The per-op file
    shrinks from \(\sim 150\) lines to \(\sim 30\).
  \item \textbf{Backend uniformity.} A single point of dispatch
    means a single place to add cross-cutting concerns: AMP
    casting, deterministic execution, profiling hooks. We added
    profiler instrumentation once, in the kernel layer, and every
    op got it.
  \item \textbf{Testability.} The kernel templates can be tested
    against synthetic scalar functions without invoking any real
    op. The unit tests in \code{test/kernel/} exercise the
    templates with mock functions and predetermined inputs.
  \item \textbf{Compile-time amortisation.} Each kernel template
    is instantiated once per arity class and once per backend, not
    once per op. Total compile time for the \code{ops/} directory
    is roughly a quarter of what it would be without the layer.
\end{itemize}

\section{The IKernel Interface}
\label{sec:kernel_abstraction:ikernel}

\code{IKernel}, declared in \code{lucid/\_C/kernel/IKernel.h}, is
the abstract base for all kernel templates.
\figref{fig:kernel_abstraction:hierarchy} sketches the hierarchy:
five arity-specialised templates above \code{IKernel} for the
regular case, plus a sibling \code{primitives/} directory of
hand-written kernels for the data-dependent operations.

\begin{figure}[t]
\centering
\begin{adjustbox}{max width=\textwidth, center}
\begin{tikzpicture}[
  base/.style={draw, rounded corners=4pt, font=\small\sffamily\bfseries,
               fill=lucid.note.bg, minimum width=3.4cm,
               minimum height=0.9cm, align=center},
  tmpl/.style={draw, rounded corners=2pt, font=\footnotesize\sffamily,
               fill=lucid.info.bg, minimum width=3cm,
               minimum height=0.7cm, align=center, inner sep=3pt},
  prim/.style={draw, rounded corners=2pt, font=\footnotesize\sffamily,
               fill=lucid.ok.bg, minimum width=3cm,
               minimum height=0.7cm, align=center, inner sep=3pt},
  inherit/.style={->, >=stealth, thick, color=lucid.accent.dark,
                  shorten >=2pt, shorten <=2pt},
]
  \node[base] (ikernel) at (0, 4.5) {IKernel};

  % Row 1: templated arity classes
  \node[tmpl] (unary) at (-6.4, 2.5) {UnaryKernel$<$Op$>$};
  \node[tmpl] (binary) at (-3.2, 2.5) {BinaryKernel$<$Op$>$};
  \node[tmpl] (reduce) at (0, 2.5) {ReduceKernel$<$Op$>$};
  \node[tmpl] (nary)   at (3.2, 2.5) {NaryKernel$<$Op,N$>$};
  \node[tmpl] (variad) at (6.4, 2.5) {VariadicKernel$<$Op$>$};

  % Row 2: data-dependent primitives
  \node[prim] (gather)  at (-6.4, 0.5) {Gather};
  \node[prim] (scatter) at (-3.2, 0.5) {Scatter};
  \node[prim] (im2col)  at (0, 0.5) {Im2Col};
  \node[prim] (sort)    at (3.2, 0.5) {Sort, TopK};
  \node[prim] (nonzero) at (6.4, 0.5) {NonZero, Unique};

  \draw[inherit] (unary) -- (ikernel);
  \draw[inherit] (binary) -- (ikernel);
  \draw[inherit] (reduce) -- (ikernel);
  \draw[inherit] (nary) -- (ikernel);
  \draw[inherit] (variad) -- (ikernel);

  \draw[inherit] (gather)  to[bend left=8]  (ikernel);
  \draw[inherit] (scatter) to[bend left=4]  (ikernel);
  \draw[inherit] (im2col)  --                (ikernel);
  \draw[inherit] (sort)    to[bend right=4] (ikernel);
  \draw[inherit] (nonzero) to[bend right=8] (ikernel);

  % Row labels (placed beneath each row, not to the side)
  \node[font=\footnotesize\itshape, color=lucid.muted, align=center]
    at (-9, 2.5) {Templated\\arity classes};
  \node[font=\footnotesize\itshape, color=lucid.muted, align=center]
    at (-9, 0.5) {Hand-written\\data-dep primitives};
\end{tikzpicture}
\end{adjustbox}
\caption[Kernel hierarchy.]{The \code{IKernel} hierarchy. The
five arity-class templates (top row) handle the regular cases ---
each unary, binary, reduction, fixed-arity, and variadic op is one
template instantiation. The data-dependent primitives (bottom row)
are full implementations rather than templates, because their
output shape or layout depends on input values; they live in
\texttt{kernel/primitives/}.}
\label{fig:kernel_abstraction:hierarchy}
\end{figure}

It is intentionally minimal:

\begin{lstlisting}[style=lucidcpp,language=C++]
class IKernel {
public:
  virtual ~IKernel() = default;

  // Allocate the output tensor(s) given input shape/dtype.
  virtual std::vector<TensorImpl> allocate_outputs(
      const std::vector<const TensorImpl*>& inputs) const = 0;

  // Run the kernel on the given inputs, writing into outputs.
  virtual void apply(
      const std::vector<const TensorImpl*>& inputs,
      const std::vector<TensorImpl*>& outputs) const = 0;
};
\end{lstlisting}

The interface deliberately abstracts away the arity: a unary
kernel sees one input and one output; a binary kernel sees two
inputs and one output; a reduction kernel sees one input, an axis
specification, and one output. Concrete kernel templates fix the
arity and the dispatch pattern.

\section{UnaryKernel}
\label{sec:kernel_abstraction:unary}

\code{UnaryKernel<Op>}, in \code{lucid/\_C/kernel/UnaryKernel.h},
is the template for single-input pointwise operations. It is
parameterised by an \code{Op} struct that defines the scalar
function and its backward:

\begin{lstlisting}[style=lucidcpp,language=C++]
struct ReluOp {
  static constexpr const char* name = "relu";

  template <typename T>
  static T forward(T x) { return x > T(0) ? x : T(0); }

  template <typename T>
  static T backward(T x, T grad_out) {
      return x > T(0) ? grad_out : T(0);
  }
};

using ReluKernel = UnaryKernel<ReluOp>;
\end{lstlisting}

\subsection{Forward dispatch}
\label{ssec:kernel_abstraction:unary_forward}

The template's forward dispatch:

\begin{enumerate}
  \item Allocate output tensor with the same shape and dtype as
    the input (with AMP override if applicable).
  \item Look up the input's device and route to the appropriate
    backend.
  \item Call the backend's \code{unary(OpKind, input)} where
    \code{OpKind} is derived from \code{Op::name}.
\end{enumerate}

The backend's \code{unary} method dispatches to the actual
implementation: for CPU, a vDSP / vForce primitive; for GPU, an
\MLX\ primitive. The template itself does not write the loop --- it
delegates to the backend, which delegates to the underlying
library's kernel.

\subsection{Backward dispatch}
\label{ssec:kernel_abstraction:unary_backward}

The backward dispatch constructs a \code{grad\_fn} node that, on
\code{backward()}, calls \code{Op::backward(saved\_input,
grad\_output)} elementwise. The saved-input list is one tensor (the
forward input). The backward kernel runs through the same template
machinery, dispatched to the same backend.

For ops where the backward depends only on the output (e.g.\
\code{tanh}, where \code{tanh'(x) = 1 - tanh(x)\textasciicircum{}2}), the
\code{Op} struct sets a flag and the autograd engine saves the
output rather than the input. This is a memory optimisation for
the common case where the output is the same size as the input.

\section{BinaryKernel}
\label{sec:kernel_abstraction:binary}

\code{BinaryKernel<Op>}, in \code{lucid/\_C/kernel/BinaryKernel.h},
is the template for two-input pointwise operations:

\begin{lstlisting}[style=lucidcpp,language=C++]
struct AddOp {
  static constexpr const char* name = "add";

  template <typename T>
  static T forward(T a, T b) { return a + b; }

  template <typename T>
  static std::pair<T,T> backward(T a, T b, T grad_out) {
      return {grad_out, grad_out};
  }
};
\end{lstlisting}

The template handles three concerns beyond what \code{UnaryKernel}
handles:

\begin{itemize}
  \item \textbf{Broadcasting.} If the two inputs have different
    shapes that are NumPy-broadcast-compatible, the kernel
    computes the broadcast shape and runs over it with strided
    access.
  \item \textbf{Type promotion.} If the two inputs have different
    dtypes, the kernel promotes both to a common dtype before
    calling the backend (\chref{ch:backend_dispatch}~Section~\ref{sec:backend_dispatch:dtype_promo}).
  \item \textbf{Cross-device check.} The kernel raises
    \code{LucidError} if the two inputs are on different devices.
\end{itemize}

The backward returns a pair of gradients, one per input. The
autograd engine threads each gradient to the corresponding input
tensor's \code{grad}.

\subsection{Broadcast unbroadcast}
\label{ssec:kernel_abstraction:unbroadcast}

A subtle case: if the inputs were broadcast in forward, the
backward must sum the gradient over the broadcast dimensions to
match each input's original shape. \code{BinaryKernel} handles
this internally: it records the broadcast pattern at forward time
and applies the reverse sum at backward time.

The handling is the same for every binary op, which is why it lives
in the template rather than in each op's file.

\section{ReduceKernel}
\label{sec:kernel_abstraction:reduce}

\code{ReduceKernel<Op>}, in \code{lucid/\_C/kernel/ReduceKernel.h},
is the template for reductions:

\begin{lstlisting}[style=lucidcpp,language=C++]
struct SumOp {
  static constexpr const char* name = "sum";

  template <typename T> static T init()         { return T(0); }
  template <typename T> static T combine(T a, T b) { return a + b; }
  template <typename T> static T finalize(T s, size_t n) { return s; }
  // (mean's finalize divides; std's finalize takes sqrt; etc.)
};
\end{lstlisting}

The reduction is parameterised by \code{init} (the identity),
\code{combine} (the binary reducer), and \code{finalize} (any
post-processing). The template handles axis selection,
\code{keepdims}, and the choice between a single-axis and
multi-axis reduction.

\subsection{Axis selection}
\label{ssec:kernel_abstraction:axis_select}

Reductions over different axes have different access patterns: a
contiguous axis reduces with unit-stride loads; a non-contiguous
axis requires strided access. The kernel inspects the input's
strides and chooses between two implementations:

\begin{itemize}
  \item Contiguous reduction: pass through to the backend's
    one-shot reduction primitive.
  \item Strided reduction: loop over outer axes manually, with
    one-shot reduction over the innermost axis.
\end{itemize}

The backend itself often makes finer choices below this level
(vDSP's \code{vDSP\_sve}-style sum has its own strided variants),
but the kernel's coarse choice keeps the dispatch tractable.

\subsection{Numerical stability}
\label{ssec:kernel_abstraction:reduce_stability}

Naive summation accumulates rounding error proportional to the
number of elements summed. For reductions over very large tensors
in FP16 or BF16, this matters: the accumulated error can exceed
the meaningful precision of the result.

The kernel handles this with a \code{ForceFloat32} promotion (the
\code{AmpPolicy} discussed in \chref{ch:op_registry}): reductions
in FP16/BF16 accumulate in FP32 and cast back at the end. The cost
is one extra cast per kernel; the benefit is correct results.

\code{logsumexp} specifically requires the standard numerical
trick (\code{logsumexp(x) = max(x) + log(sum(exp(x - max(x))))}),
which the kernel implements directly rather than chaining the
naive composition.

\section{NaryKernel and VariadicKernel}
\label{sec:kernel_abstraction:variadic}

The remaining two kernel templates handle ops with more than two
inputs:

\begin{itemize}
  \item \code{NaryKernel<Op, N>}: fixed arity \(N\). Used for
    operations like \code{addcmul} (\(N = 3\), three same-shape
    pointwise inputs) and \code{where} (\(N = 3\), condition plus
    two same-shape value tensors).
  \item \code{VariadicKernel<Op>}: variable arity. Used for
    operations like \code{cat} (concatenate any number of tensors
    along an axis) and \code{stack} (stack any number of tensors
    into a new axis).
\end{itemize}

The two templates share most of their implementation. The split
exists because the fixed-arity case can use stack-allocated input
pointers, while the variable-arity case needs a vector.

\section{Contiguity Inference}
\label{sec:kernel_abstraction:contig}

\code{Contig.h} in the kernel directory provides the contiguity
queries that the kernel templates consult before deciding between
a fast path (assume contiguous, call backend's one-shot kernel)
and a slow path (handle strided access manually).

\subsection{The is\_contig\_along predicate}
\label{ssec:kernel_abstraction:is_contig_along}

The query is more nuanced than ``is the whole tensor contiguous.''
For a binary kernel, what matters is whether \emph{both inputs} are
contiguous along the axis the kernel iterates over. For a reduction
kernel, what matters is whether the reduced axis is contiguous.
\code{Contig.h} provides predicates for these specific cases:

\begin{lstlisting}[style=lucidcpp,language=C++]
bool is_contiguous(const TensorImpl&);
bool is_contiguous_along(const TensorImpl&, int axis);
bool is_unit_stride_innermost(const TensorImpl&);
bool both_contiguous(const TensorImpl&, const TensorImpl&);
\end{lstlisting}

The predicates are constant-time after the first call (the result
is cached on the \code{TensorImpl}).

\subsection{When to materialise}
\label{ssec:kernel_abstraction:materialise}

If the kernel cannot run on a non-contiguous view, it materialises
the input via the backend's \code{contiguous} primitive. The cost
is one copy; the alternative would be a slow kernel that handles
arbitrary strides, which is generally not worth the implementation
complexity.

The materialise-vs-strided decision is made per-kernel, not
per-op. Some kernel templates accept strided input (most
elementwise); others always materialise (the convolution backward,
the FFT). The choice is documented in the kernel template's
implementation.

\section{Primitive Kernels (Data-Dependent Operations)}
\label{sec:kernel_abstraction:primitives}

Not every operation fits the templated kernel pattern.
Data-dependent operations --- where the output shape or layout
depends on the input values, not just the input shapes --- have
their own files in \code{lucid/\_C/kernel/primitives/}. The list:

\begin{itemize}
  \item \code{Gather}: \code{out[i, j] = input[i, indices[i, j]]}.
    Output shape comes from the indices, not from the input.
  \item \code{Scatter}: the inverse; writes \code{updates} into
    \code{input} at positions given by \code{indices}.
  \item \code{Im2Col}: rearranges convolution inputs into a column
    matrix for GEMM-based convolution.
  \item \code{Pad}: adds padding around a tensor.
  \item \code{Sort}: sorts along an axis; output values are
    permuted.
  \item \code{NonZero}: returns the coordinates of nonzero
    entries; output shape depends on the input values.
  \item \code{Unique}: returns the unique values in a tensor;
    output shape is at most input shape.
  \item \code{SearchSorted}: binary-searches one tensor in
    another.
\end{itemize}

These primitives are full implementations, not templates: each one
has its own forward and backward code, its own dispatch to the
backend's specialised functions, and (in the case of \code{NonZero}
and \code{Unique}) its own CPU-stream forcing because the output
shape requires reading the input values.

\subsection{The CPU stream forcing}
\label{ssec:kernel_abstraction:cpu_stream_force}

For \code{NonZero} and \code{Unique} on metal tensors, the
dispatcher's data-dependent-output carve-out
(\chref{ch:backend_dispatch}~Section~\ref{ssec:backend_dispatch:data_dep_carveout})
bridges the tensor to CPU, runs the primitive, and bridges the
result back. The kernel itself sees only CPU tensors; the bridging
is the dispatcher's responsibility.

\section{Summary and Forward References}
\label{sec:kernel_abstraction:summary}

The kernel abstraction layer extracts the common control structure
from \Lucid's operations into a small set of templates:
\code{UnaryKernel}, \code{BinaryKernel}, \code{ReduceKernel},
\code{NaryKernel}, \code{VariadicKernel}. Each template
parameterises over an \code{Op} struct that defines the scalar
function and (where applicable) its backward. Contiguity inference
keeps the templates from running slow paths on contiguous inputs,
and a small set of data-dependent primitives lives outside the
templated machinery for the cases that do not fit.

The next four chapters describe the op categories that use these
templates: \chref{ch:unary_functions} for unary ops,
\chref{ch:binary_functions} for binary ops, \chref{ch:reductions}
for reductions, and \chref{ch:shape_view} for shape and broadcast
ops. \chref{ch:indexing} treats the indexing primitives, and
\chref{ch:composite_ops} treats the pure-Python composite layer.

\begin{lucidnote}
For the user: the kernel layer is invisible. The user calls
\code{lucid.relu(x)}; the call passes through the dispatcher to
\code{UnaryKernel<ReluOp>}, which dispatches to the backend's
\code{relu} primitive. The whole chain runs in microseconds for
non-trivial inputs.
\end{lucidnote}
